\section*{第 \ref{ch:b3:adaptive-immunity} 章——适应性免疫与疫苗接种}

\begin{solution}{exo:b3:adaptive-immunity:1}
每个\omterm{def:b3:adaptive-immunity:lymphocytes}{淋巴细胞}只有一种特异性（诺索尔与莱德伯格：每个细胞
只分泌针对一种\omterm{def:b3:adaptive-immunity:lymphocytes}{抗原}的\omterm{def:b3:adaptive-immunity:antibody}{抗体}）；受体在遇到\omterm{def:b3:adaptive-immunity:lymphocytes}{抗原}之前已经形成（利根川：
基因是在 B 细胞中重排的，而不是响应\omterm{def:b3:adaptive-immunity:lymphocytes}{抗原}才重排）；
\omterm{def:b3:adaptive-immunity:lymphocytes}{抗原}选择并扩增相应的克隆，使其分化为效应细胞和\omterm{def:b3:adaptive-immunity:response}{记忆细胞}；
自身反应性克隆
在发育过程中被清除。
\end{solution}

\begin{solution}{exo:b3:adaptive-immunity:2}
两条重链和两条轻链，每条链末端为可变结构域，其后为恒定结构域；
由六个超变环构成两个相同的结合位点；Fc 段供其他细胞和\omterm{def:b3:innate-immunity:complement}{补体}识别。
IgM：最早出现的\omterm{def:b3:adaptive-immunity:antibody}{抗体}，为五聚体，可激活\omterm{def:b3:innate-immunity:complement}{补体}。
IgG：血液和组织中的主要\omterm{def:b3:adaptive-immunity:antibody}{抗体}，具有\omterm{def:b3:innate-immunity:complement}{调理作用}和中和作用，可穿过胎盘。
IgA：以二聚体形式分泌到黏膜表面及乳汁中。
IgE：与肥大细胞结合，抵御寄生虫，是过敏反应的介导者。
IgD：存在于初始 B 细胞表面，
功能较次要，已知功能较少。
\end{solution}

\begin{solution}{exo:b3:adaptive-immunity:3}
MHC~I 类分子：所有有核细胞表达；肽段来自胞质蛋白，经蛋白酶体切割，
长 8--10 个残基；由 CD8 细胞毒性 T 细胞识别。MHC~II 类分子：\omterm{def:b3:innate-immunity:innate}{树突状细胞}、\omterm{def:b3:innate-immunity:innate}{巨噬细胞}、
B 细胞表达；肽段来自摄入并在\omterm{def:b3:membrane-traffic:endocytosis}{内体}中降解的蛋白，
长 13--25 个残基；由 CD4 辅助性 T 细胞识别。
\end{solution}

\begin{solution}{exo:b3:adaptive-immunity:4}
$R_{0}$ 为完全易感人群中一个病例平均所产生的病例数；
阈值为 $p_{c} = 1 - 1/R_{0}$。腮腺炎
\qty{80}{\%}；百日咳 \qty{92}{\%}；1918 年流感 \qty{60}{\%}。
\end{solution}

\begin{solution}{exo:b3:adaptive-immunity:5}
$50\times 20\times 5 = 5000$ 条重链，$60\times 4 = 240$ 条轻链：
组合数为 $1.2\times 10^{6}$；计入\omterm{def:b3:adaptive-immunity:antibody}{连接多样性}后达 $1.2\times
10^{10}$，为 B 细胞数的十倍，因此动物至多只能表达其可产生的
组合的十分之一，几乎每个细胞
都带有独特的受体。
\end{solution}

\begin{solution}{exo:b3:adaptive-immunity:6}
$1000\times 10^{-3} = 1$ 次突变（每个细胞每次分裂），二十次分裂后为 $20$ 次。
在 $10^{5}$ 个细胞中，$2\times 10^{6}$ 次突变事件落在
$3000$ 种可能的单一替换上：每种替换都被取样数百次，
另有双重与三重组合——\omterm{def:b3:adaptive-immunity:germinal-centre}{生发中心}
对受体邻域进行了穷尽式探索。
\end{solution}

\begin{solution}{exo:b3:adaptive-immunity:7}
T 细胞在\omterm{def:b3:adaptive-immunity:tolerance}{胸腺}中被选择，以识别宿主自身 MHC 等位基因上
呈递的肽段；他人的等位基因具有形状不同的凹槽，
能以不同方式容纳不同的肽段，因此同一\omterm{def:b3:virology:virus}{病毒}肽段
不会被识别。移植物的外来 MHC 分子本身会被大量
T 细胞识别为“携带异常肽段的自身 MHC”，因此引发排斥。
这种多态性之所以得以维持，是因为具有多种等位基因的种群
能呈递任何病原体肽段的更广泛样本，且没有任何病原体能逃避所有人；
杂合子呈递的范围为纯合子的两倍，
因而受到选择的青睐。
\end{solution}

\begin{solution}{exo:b3:adaptive-immunity:8}
(a) 缺乏 V(D)J 重组：没有 B 细胞或 T 细胞，导致严重联合免疫缺陷；
(b) 缺乏 AIRE：\omterm{def:b3:adaptive-immunity:tolerance}{胸腺}细胞无法展示组织\omterm{def:b3:adaptive-immunity:lymphocytes}{抗原}，自身反应性 T 细胞逃逸，
引发针对多种器官的\omterm{def:b3:adaptive-immunity:tolerance}{自身免疫}；
(c) 缺乏 \omterm{def:b3:adaptive-immunity:tolerance}{FoxP3}：没有调节性 T 细胞，导致早发且致命的\omterm{def:b3:adaptive-immunity:tolerance}{自身免疫}和过敏；
(d) 缺乏 AID：无法发生\omterm{def:b3:adaptive-immunity:antibody}{类别转换}和超突变——IgM 大量存在，
几乎没有 IgG 或 IgA，\omterm{def:b3:adaptive-immunity:antibody}{抗体}亲和力低，反复感染；
(e) 缺乏 CD4 T 细胞：无法为 B 细胞或\omterm{def:b3:innate-immunity:innate}{巨噬细胞}提供辅助，
\omterm{def:b3:adaptive-immunity:antibody}{抗体}应答不良，易受机会性感染，
即艾滋病的免疫缺陷。
\end{solution}

\begin{solution}{exo:b3:adaptive-immunity:9}
$p_{c} = 1 - 1/6 = 0.83$；覆盖度为 $0.83/0.9 = 93\,\%$。当覆盖度为 \qty{80}{\%}
时，免疫个体比例为 $0.72$，$R_{\text{有效}} = 6\times 0.28
= 1.7$。在易感者中，疫情的传播如同 $R_{0}$ 为
$1.7$ 一样（免疫者只是稀释了接触），并由终值
关系 $s_{\infty} - \ln s_{\infty}/1.7 = 1$ 可得 $s_{\infty}
\approx 0.31$：约 \qty{69}{\%} 的易感者、
\qty{19}{\%} 的总人口被感染。
\end{solution}

\begin{solution}{exo:b3:adaptive-immunity:10}
由 $\mathrm{d}s/\mathrm{d}t = -\beta si$ 与 $\mathrm{d}i/\mathrm{d}t
= \beta si - \gamma i$，得 $\mathrm{d}i/\mathrm{d}s = -1 + \gamma/(\beta s)$，
故 $i + s - (\ln s)/R_{0}$ 守恒；取 $s = 1$、$i = 0$ 为起点，
终点处 $i = 0$、$s = s_{\infty}$，于是 $s_{\infty} - \ln
s_{\infty}/R_{0} = 1$。解得 $R_{0} = 1.5$ 时，$s_{\infty} = 0.42$
（\qty{58}{\%} 被感染）；$2$ 时 $0.20$（\qty{80}{\%}）；$3$ 时 $0.06$
（\qty{94}{\%}）；$5$ 时 $0.007$（\qty{99.3}{\%}）。并非所有人都会被感染：
随着易感者被耗尽，有效再生数在仍有部分易感者时降至 1 以下，
疫情在波及他们之前就已消退。
\end{solution}

\begin{solution}{exo:b3:adaptive-immunity:11}
$1.5^{n} = 1000$：$n = \ln 1000/\ln 1.5 \approx 17$，
即需连续 17 次改进。若每个细胞每次分裂产生一次突变，而每一百次突变中
仅有一次能够改进，则每轮至少需要筛选一百个细胞才能找到一个；
实际上需要数千个，因为一半的突变会破坏结合，且选择过程带有噪声。
17 轮，每轮两次分裂，约需九天——
考虑到各种低效因素后，与观察到的两至三周相吻合。
循环将突变（暗区）与检测（明区）分开，
使每个变体在被允许增殖之前先接受\omterm{def:b3:adaptive-immunity:lymphocytes}{抗原}的检验，
恰如遗传算法交替进行
变异与选择。
\end{solution}

\begin{solution}{exo:b3:adaptive-immunity:12}
女性：若干免疫基因位于 X 染色体上，其中一些（TLR7）
可逃避失活，两份拷贝均能表达，且雌激素
增强 B~细胞的存活——检验：多一条 X 染色体的男性（克兰费尔特综合征）
患狼疮的风险与女性相似，事实确实如此。\omterm{def:b3:adaptive-immunity:mhc}{HLA}：相关的
等位基因能很好地呈递相关自身肽，或在\omterm{def:b3:adaptive-immunity:tolerance}{胸腺}中未能呈递
它们，从而使相应克隆不被清除——检验：小鼠
携带人类该等位基因的\omterm{def:b3:genetic-engineering:transgenic}{转基因}小鼠，在接受
\omterm{def:b3:adaptive-immunity:lymphocytes}{抗原}后发病。发病率上升，原因在于早期微生物接触减少，以及
\omterm{def:b3:microbiomes:symbiosis}{微生物组}改变（\omterm{def:b3:bacteriology:antibiotics}{抗生素}、饮食、剖宫产）使调节性
T~细胞减少、\omterm{def:b3:innate-immunity:inflammation}{炎症}增多——检验：
无菌或使用\omterm{def:b3:bacteriology:antibiotics}{抗生素}饲养的易患糖尿病小鼠更多发病，移民的
子女在一代之内即
获得新居住国的发病率。
\end{solution}

\begin{solution}{pb:b3:adaptive-immunity:1}
\textbf{1.} $6000\times 320 = 1.9\times 10^{6}$；$\times 3\times 10^{4}
= 5.8\times 10^{10}$。
\textbf{2.} $10^{11}$ 个细胞与 $5.8\times 10^{10}$ 种可能相比：一个机体
足以容纳每种受体约两次，但各克隆数量不一，多数
序列恰好从未被产生，因此每个机体表达的是大而
不完整的样本——两人之间共有的
受体仅占很小比例。
\textbf{3.} $10^{11}/10^{5} = 10^{6}$ 个细胞；在 $10^{8}$
个细胞的\omterm{def:b3:adaptive-immunity:lymphocytes}{淋巴结}中约 $1000$ 个。
\textbf{4.} $10^{9}/10^{5} = 10^{4}$ 个前体细胞——少一百倍；
新生儿的滤泡和\omterm{def:b3:innate-immunity:innate}{树突状细胞}也尚未成熟，
缺乏记忆，T~细胞的辅助有限，因此应答缓慢且弱，
母源 IgG 弥补了这一缺口。
\textbf{5.} 它横跨 V--D 与 D--J 连接处，此处核苷酸
被随机删除或添加：它的序列在基因组中
并无编码。它位于两条链之间，居于结合位点中心，
是最常与表位接触的环。
\textbf{6.} \omterm{def:b3:adaptive-immunity:lymphocytes}{抗原}出现之前，多样性必须足够广，使得
无论来者为何，总有受体能与之结合；\omterm{def:b3:adaptive-immunity:lymphocytes}{抗原}出现之后，
更宜把精力用于改造已知能结合的少数受体。先粗后精，
远比单独采用其中任何一种高效。
\textbf{7.} 七天即 $21$ 次分裂：$100\times 2^{21} = 2.1\times
10^{8}$ 个细胞；\qty{30}{\%}，$6.3\times 10^{7}$ 个\omterm{def:b3:adaptive-immunity:response}{浆细胞}。
\textbf{8.} $6.3\times 10^{7}\times 2000\times \num{86400} = 1.1\times
10^{16}$ 个分子每天；$\times 1.5\times 10^{5}\times 1.66\times
10^{-24}$ 克 $= \qty{2.7}{mg}$；在 \qty{3}{L} 中，每天 \qty{0.9}{mg/L}。
\textbf{9.} \qty{27}{mg} 在 \qty{3}{L} 中：\qty{9}{mg/L} $= \qty{0.009}{g/L}$，
仅占总 IgG 的千分之一——针对任一\omterm{def:b3:adaptive-immunity:lymphocytes}{抗原}的\omterm{def:b3:adaptive-immunity:antibody}{抗体}
只占整体中很小的一部分。
\textbf{10.} 百倍即 $\log_{2}100 = 6.6$ 个半衰期，约
$140$ 天。持续数年的\omterm{def:b3:adaptive-immunity:antibody}{抗体}来自骨髓中的
长寿\omterm{def:b3:adaptive-immunity:response}{浆细胞}（持续分泌）以及
在再次接触时被再刺激的\omterm{def:b3:adaptive-immunity:response}{记忆细胞}。
\textbf{11.} 九次分裂：$10^{5}\times 512 = 5\times 10^{7}$ 个细胞，
即第 3 天的数量，是初次应答第 3 天 $5\times 10^{4}$ 的一千倍，
也只是初次应答第 7 天水平的四分之一。
\textbf{12.} 重排后的 VDJ 紧邻 C$_{\mu}$，因此最初的
转录本为 IgM；转换为 C$_{\gamma}$ 需要 T~细胞辅助及 AID，
这一过程需要数天。\omterm{def:b3:adaptive-immunity:response}{记忆细胞}已完成转换，因此其子代
可立即产生 IgG。
\textbf{13.} $1000\times 10^{-3} = 1$，即每次分裂一次突变；二十次分裂共 $20$ 次。
\textbf{14.} 百分之一；$10^{4}\times 0.01 = 100$ 个
改进细胞每次分裂产生。
\textbf{15.} $1.5^{n} = 1000$：$n \approx 17$；$17\times \qty{12}{h}
\approx 8.5$ 天。
\textbf{16.} 一半的子代丧失结合能力：若无选择，$0.5^{20}
\approx 10^{-6}$ 的谱系在二十次分裂后仍能结合。
明区必须在每个循环中淘汰失败者，
使得只有功能正常且经改进的受体得以延续。
\textbf{17.} 第一剂之后留下的记忆 B~细胞亲和力已提高，且已完成
\omterm{def:b3:adaptive-immunity:antibody}{类别转换}；第二剂将它们带入新的\omterm{def:b3:adaptive-immunity:germinal-centre}{生发中心}，
\omterm{def:b3:adaptive-immunity:germinal-centre}{体细胞高频突变}与选择从更高的起点重新开始，
因此随后产生的\omterm{def:b3:adaptive-immunity:antibody}{抗体}更加优良。
\textbf{18.} 部分感染者最终产生了针对
\omterm{def:b3:virology:virus}{病毒}无法改变的保守位点的\omterm{def:b3:adaptive-immunity:antibody}{抗体}（广谱中和\omterm{def:b3:adaptive-immunity:antibody}{抗体}），
但需经数年的成熟。\omterm{def:b3:adaptive-immunity:germinal-centre}{生发中心}可以被
引导：一系列免疫原首先吸引正确的种系前体，
随后奖励与保守位点结合者，即可沿此路径引导成熟，
数月而非数年即可完成。
\textbf{19.} $p_{c} = 1 - 1/15 = 93.3\,\%$。一剂：$0.933/0.93 =
100.4\,\%$——无法达到；两剂：$0.933/0.97 = 96.2\,\%$。
\textbf{20.} 免疫比例 $0.90\times 0.97 = 0.873$；$R_{\text{eff}} = 15\times
0.127 = 1.9$：麻疹仍在传播。
\textbf{21.} $60\,000\times (0.10 + 0.90\times 0.03) = 7600$ 名易感者
每个出生队列；五年后为 $38\,000$——一个储备库，
输入性病例可在其接触者中找到足够的易感者，
以维持传播链。
\textbf{22.} 在易感者中，$s_{\infty} - \ln s_{\infty}/1.9 =
1$ 得 $s_{\infty} \approx 0.23$：约占其中的 \qty{77}{\%}，
即约 $29\,000$ 例病例。
\textbf{23.} 免疫比例 $0.95\times 0.97 = 0.92$；$R_{\text{eff}} = 15\times
0.08 = 1.2$——仍大于 1；阈值需要两剂接种达到 \qty{96}{\%}，
而麻疹正是对最后几个百分点格外苛刻的
疾病。
\textbf{24.} 婴儿仅靠周围人的免疫力获得保护
（以及数月内的母源\omterm{def:b3:adaptive-immunity:antibody}{抗体}）：若传播链无法形成，
病例便无法传到他们身上。在 \qty{85}{\%} 覆盖度下，
免疫者比例为 $0.82$，$R_{\text{eff}} = 2.6$：疫情
持续流行，而未接种\omterm{def:b3:adaptive-immunity:vaccine}{疫苗}、最为脆弱的婴儿遭到感染。
\textbf{25.} 受体库约 $6\times 10^{10}$；一次应答的\omterm{def:b3:adaptive-immunity:response}{浆细胞}
每天产生 \qty{2.7}{mg} 的 IgG；两剂接种覆盖率 \qty{96}{\%}。
\end{solution}
